% Generated from locale/tr/content/first-order-logic/sequent-calculus/rules-and-proofs.tex; do not hand-edit.


\iftag{FOL}
      {\olfileid[tr]{fol}{seq}{rul}}
      {\olfileid[tr]{pl}{seq}{rul}}

\olsection{Kurallar ve \usetoken{P}{derivation}}

Aşağıda $\Gamma, \Delta, \Pi, \Lambda$, öğeleri !!{sentence}s olan sonlu
dizileri göstersin.

\begin{defn}[Sekant]
Bir \emph{sekant},
\[
\Gamma \Sequent \Delta
\]
biçiminde bir ifadedir; burada $\Gamma$ ile $\Delta$, öğeleri $\Lang L$
dilindeki !!{sentence}s olan sonlu (boş da olabilen) dizilerdir. $\Gamma$'ya
\emph{öncül}, $\Delta$'ya ise \emph{ardıl} denir.
\end{defn}

\begin{explain}
Bir sekantın ardındaki sezgisel düşünce şudur: öncüldeki bütün !!{sentence}s
doğruysa ardıldaki !!{sentence}s arasından en az biri doğrudur. Başka bir
deyişle $\Gamma = \tuple{!A_1, \dots, !A_m}$ ve $\Delta = \tuple{!B_1,
\dots, !B_n}$ ise $\Gamma \Sequent \Delta$ ancak ve ancak
\[
(!A_1 \land \cdots \land !A_m) \lif (!B_1 \lor \cdots \lor
!B_n)
\]
doğru olduğunda doğrudur. İki özel durum vardır: $\Gamma$'nın boş olması ve
$\Delta$'nın boş olması. $\Gamma$ boşsa, yani $m=0$ ise $\quad \Sequent
\Delta$ ancak ve ancak $!B_1 \lor \dots \lor !B_n$ doğru olduğunda
doğrudur. $\Delta$ boşsa, yani $n=0$ ise $\Gamma \Sequent \quad$ ancak ve
ancak $\lnot(!A_1 \land \dots \land !A_m)$ doğru olduğunda doğrudur.
Karşılık gelen !!{sentence} geçerliyse sekanta geçerli deriz.
\end{explain}

$\Gamma$, öğeleri !!{sentence}s olan bir diziyse $\Gamma,!A$ ile $\Gamma$'nın sağ ucuna
$!A$ eklenerek elde edilen diziyi (ve $!A,\Gamma$ ile $\Gamma$'nın sol ucuna
$!A$ eklenerek elde edilen diziyi) gösteririz. $\Delta$ da öğeleri !!{sentence}s olan bir
dizisiyse $\Gamma,\Delta$, iki dizinin art arda eklenmesidir.

\begin{defn}[Başlangıç Sekantı]
Bir \emph{başlangıç sekantı}, herhangi bir !!{sentence}~$!A$ için
\iftag{prvFalse,prvTrue}{aşağıdaki biçimlerden birindeki sekanttır:
  \begin{enumerate}
    \item $!A \Sequent !A$
    \tagitem{prvTrue}{$\quad \Sequent \ltrue$}{}
    \tagitem{prvFalse}{$\lfalse \Sequent \quad$}{}
  \end{enumerate}}
{$!A \Sequent !A$ biçimindeki sekanttır}.
\end{defn}

Sekant hesabındaki !!^{derivation}s, en üstteki sekantların başlangıç
sekantları olduğu belirli sekant ağaçlarıdır; bir sekant bir veya iki başka
sekantın altında duruyorsa onlardan bir çıkarım kuralıyla doğru biçimde elde
edilmelidir. $\Log{LK}$ kuralları iki ana türe ayrılır: \emph{mantıksal}
kurallar ve \emph{yapısal} kurallar. Mantıksal kurallar, alttaki sekantta
!!{main operator} konumunda duran ve $!A$ ve/veya $!B$ içeren
!!{sentence} temel alınarak adlandırılır. Her kuralın iki biçimi vardır: biri
!!{sentence} içindeki !!{operator} sol tarafta bulunan bir sekant, diğeri bu !!{sentence}
sağ tarafta bulunan bir sekant çıkarmak içindir.

